\begin{table}[htbp]
\centering\small
\caption{Coverage of recorded safety documentation by follow-up window.}\label{tab:safety_documentation}
\begin{tabular}{p{31mm}rrrrr}\toprule
Initial treatment & Initial N & \shortstack{With note through\\W24, n/N} & \shortstack{No note through\\W24, n} & \shortstack{With note at any\\follow-up, n/N} & \shortstack{Note records\\through W24, n} \\ \midrule
Dupilumab & 65 & 27/65 & 38 & 28/65 & 30 \\
Abrocitinib 100 mg & 56 & 20/56 & 36 & 20/56 & 28 \\
Abrocitinib 200 mg & 42 & 19/42 & 23 & 19/42 & 37 \\
\bottomrule\end{tabular}
\par\small\raggedright Patients are counted once within each documentation window. The through-W24 window uses recorded nominal visits up to W24; the all-follow-up column also includes later recorded visits. Note records can repeat within a patient and are not adverse-event counts. Absence of a note means unknown documentation, not absence of an event. Initial-group N is a cohort denominator, not a verified safety-observation denominator. These counts do not estimate AE or SAE incidence, severity or causality, and do not establish event-time treatment attribution. Source, \texttt{safety\_documentation\_windows.csv}.
\end{table}
